\section{Verification for the five complete schedules}
\label{sec:schedules}

The C1--C4 applications below have written source proofs; C5's clock and
regularity are also formalized for its declarative Lean schedule. No
application is inferred from numerical tests. We use the pinned
\texttt{schedules.rs} constructors and
\texttt{model.rs} movement convention. A positioning call takes the maximum
absolute displacement of the participating heads; a head that arrives
early stays at its target. Every tick is charged, including ticks whose
annotated data operation is a no-op.

Set
\[
 K=m-1,\quad M_r=L/2^r\ (1\le r\le K),\quad
 M_K=2db\ge8,\quad \sum_{r=1}^K M_r=L-2db.
\]
The root owns preparation and final positioning/output. Depth $r$ owns
the complete forward block $F_r$ and backward block $B_r$, including their
entry moves and internal wraps. Deeper stages are empty. These intervals
partition the actual timeline. Take $V_r=L/2^r$ at every depth, even after
the last nonempty stage. Within a depth-$r$ stage, abbreviate $M=M_r$.

\begin{table}[htbp]
\centering\small
\begin{tabular}{lrrrr}
\toprule
Constructor & tapes & legal dyadic $d$ & $|F_r|+|B_r|$ & $a$ \\
\midrule
\texttt{c1\_torus\_cascade} & 2 & $\ge16$ & $\le128M_r$ & 128 \\
\texttt{c2\_bfly\_cascade} & 2 & $\ge4$ & $\le128M_r$ & 128 \\
\texttt{c3\_ten\_entry\_cascade} & 3 & $\ge4$ & $\le80M_r$ & 80 \\
\texttt{c4\_no\_overwrite\_cascade} & 3 & $\ge4$ & $\le128M_r$ & 128 \\
\texttt{c5\_comb\_fanout} & 2 & $\ge4$ & $\le18M_r$ & 18 \\
\bottomrule
\end{tabular}
\caption{Safe all-tick budgets. Root work is also at most $aL$.
All geometric bounds hold for integer $b\ge1$. The legal sizes are
\emph{not} numerical exclusion thresholds.}
\label{tab:bounds}
\end{table}

\subsection{The common variable-speed loop}
\label{sec:variable}

The forward loops in C1--C4 use counters $c_0,c_1$, with cap $M-1$.
Put $u=\lfloor M/3\rfloor$. At local pre-step time $t$, their exact
positions relative to the starting cells are
\begin{equation}
\label{eq:variable}
 c_0(t)=\min(t,M-1),\qquad
 c_1(t)=\min\bigl(M-1,\max(\lceil t/2\rceil,t-u+1)\bigr).
\end{equation}
For C1--C3 the loop has $M+u-2$ transitions and this formula includes
its final configuration. C4 adds one terminal stationary event, giving
$M+u-1$ transitions.

Here is a direct verification of the source recurrence. Maintain
$c_1\le c_0$. While $c_0\le M-2$, write $M=3u+v$, $0\le v\le2$.
Then $c_0\le3u$ and
\[
 \lfloor2c_0/3\rfloor+u-c_0
   =u-\lceil c_0/3\rceil\ge0.
\]
The source's condition for advancing head~0 is therefore true whenever
it is unfinished. Before $c_1$ reaches $u$, head~1 advances precisely
when $c_0\ge2c_1$; its successive values are $\lceil t/2\rceil$.
At $t=2u-1$ it reaches $u$, then advances every tick until capped,
giving $t-u+1$. The caps preserve $c_1\le c_0$, and some unfinished
head always moves, so the fallback is unused. Both caps are first
reached at $t=M+u-2$. This proves the formula and endpoint counts.

After scaling time and space by $M$, the frozen profiles converge
uniformly, with error $O(1/M)$, to
\begin{equation}
\label{eq:variable-limit}
 x(s)=\min(s,1),\qquad y(s)=\min(1,\max(s/2,s-1/3)).
\end{equation}
The breaks are $2/3,1,4/3$. Monotonicity alone would not establish this
finite-affine limit; the recurrence does.

\subsection{C1 and C2: preparation and the shared two-tape core}

\paragraph{C1's growing row loop.}
One torus-shift helper first rotates columns within rows. Put $w=db$,
$\ell_1=(d-v)b$, $\ell_2=vb$ for $0<v<d$. A nonfinal row costs
\[
 \ell_1+(w-1)+\ell_2+(w+1)=3w:
\]
two scans, the within-row rewind, and the positioning for the next row.
The final row costs $2w-1$. The whole phase therefore costs exactly
$3L-w-1$. At the nonfinal row-boundary times $3iw$, both heads have
advanced $iw$ relative to their initial positions. Between boundaries,
their distances from the speed-$1/3$ profile are at most $2w+2$ and
$w+2$, respectively; these bounds also cover the frozen final extension
to $3L$. Since $w/L=1/d$, both scaled paths converge uniformly to
$s\mapsto\min(s/3,1)$.

The following temporary-bank copy is a straight $L$-tick scan. The row
rotation has two straight portions of total length $L$ and a rewind of
length $L-1$. Its entry moves are paid. For $v=0$, the column rotation
is instead one straight scan. Each helper has a finite-affine limit
whenever its row-rotation fraction converges.
For dyadic $d\ge256$ the four actual parameter pairs $(u,v)$ are
\[
 (3d/16+1,16),\ (d-16,13d/16-1),\
 (7d/16+2,22),\ (d-22,9d/16-2).
\]
Their fractions $u/d$ converge. The near-full rotations retain their
macroscopic rewinds; a small complementary segment does not make the
whole helper negligible.

Helper $j\in\{1,2,3,4\}$ uses temporary bases $j(L+4)$.
Its three large positionings cost at most $3(j+1)(L+4)$; its loops
cost at most $6L$. All four helpers therefore cost at most $66L+168$.
Loading, seed positioning and seed copy bring preparation to at most
$69L+173\le70L$. The shared core starts at work positions
$(5(L+4)+L-1,L-1)$, with frame base $5(L+4)$.
The eventual parameter formula $d\ge256$ is used for convergence only;
it is not an exclusion cutoff.

\paragraph{C2's fixed prelude.}
After loading and full reversal, C2 has two shifted forward/reverse
pairs with shifts
\[
 \Delta_{\rm M}=(d/2)(d-1)b,\qquad \Delta_{\rm L}=(d-1)b.
\]
Their ratios to $L$ tend to $1/2$ and $0$. Each pair has two scans of
length $L-\Delta$. In particular, the two small-shift positionings vanish
after scaling, but their two scans each still have limiting duration~1.
Loading, reversal and the four scans cost at most $6L$; their five
positionings at most $5L$; seed positioning and copy at most $2L+4$.
Thus preparation costs at most $14L$ and has finitely many affine
limiting pieces. Core entry is $(L+4+L-1,L-1)$, with frame base $L+4$.

\paragraph{The complete shared stages.}
At depth $r$, the second tape allocates
\[
 Z_r=\mathrm{scr}_r+M_r+2,\quad
 \mathrm{comb}_r=\mathrm{scr}_r+2M_r+4,\quad
 \mathrm{scr}_{r+1}=\mathrm{scr}_r+3M_r+6.
\]
Allocating a name takes no tick, but all moves to these addresses do.
In $F_r$, tape~0 stays in a range $3M+3$. Tape~1's range is at most
$4M+6$ except at the first depth, where the C1 prelude requires the safe
bound $10M+22$ (C2 is smaller). There are at most five positioning calls
and at most $5M$ loop events. Hence
\[
 |F_r|\le55M+110.
\]
In $B_r$, tape~0's range is at most $3M+3$. At a nondeepest entry,
tape~1 is at
$\mathrm{scr}_r+4M+9+\mathrm{shift}_{r+1}$, with
$1\le\mathrm{shift}_{r+1}<M/2$; its subsequent positions are between
$Z_r$ and $\mathrm{comb}_r+M-1$. A range $4M+10$ covers the whole
stage, also at the deepest entry. At most six positionings and $8M$
loop events give
\[
 |B_r|\le32M+60,\qquad |F_r|+|B_r|\le87M+170\le128M.
\]
Here $M\ge8$. The shared routine's variable named \texttt{comb} is a
straight $M$-cell copy, not a growing tooth loop.

At core exit tape~1 is at $\mathrm{frame\_base}+L+\mathrm{shift}_1+3$;
tape~0 is at the last cell of its final source. The final positioning
and output cost at most $9L$ for C1 and $5L$ for C2. Root work is thus
at most $79L$ and $19L$, below $128L$ in both cases.

For each fixed $r$, the source shifts
$\mathrm{shift}_r=(\lfloor M_r/3\rfloor+b)\bmod M_r$ have ratio
$\mathrm{shift}_r/M_r\to1/3$, since $b/M_r=2^r/d^2\to0$.
Normalized entry offsets and durations converge. All ordinary copies,
reversals and positionings have finite-affine limits; doubling has at most
one cell of two-tick raster error from its affine profiles.
Equation~\eqref{eq:variable-limit}
handles the remaining loop. Finite concatenation, including freezing at
the stage end, proves \eqref{eq:profile} at every fixed depth and at
the root.

\subsection{C3: bounded-entry geometry still has regular stages}

For $F_r$, write $A_0$ for the current tape-0 base, $H$ for the scratch
base and $Z$ for the tape-2 frame. The full coordinate ranges are
\[
 [A_0,A_0+3M-1],\quad[H-2M-5,H+M-1],\quad[Z,Z+M].
\]
The wider tape-1 range includes its incoming position at every depth.
At most six positionings cost $4M+6$ each. The variable loop, the two
shifted-copy pieces, their possible joining event, and copy-back cost
at most $4M+1$. Thus $|F_r|\le28M+37\le40M$.

For $B_r$, write $A_0$ for the current tape-0 base and $H$ for the comb
base. The full ranges are
\[
 [A_0,A_0+4M+1],\quad[H-M-5,H+M-1],\quad[Z,Z+M].
\]
There are $2^r$ teeth, each of length $M/2^r=d^2b/4^r\ge4b$.
The comb has $3M-2\cdot2^r-1$ ticks: its last tooth has no final join.
It returns the three heads to $(A_0,H,Z)$. Doubling and the two reversed
pieces cost $4M$ and return tape~0 to its new base $A_0+2M+2$.
At most five positionings therefore give
$|B_r|\le5(4M+6)+7M\le40M$. The last shifted reverse can leave
tape~1 at a nonzero comb offset; the next stage's range above includes it.

Preparation costs $5L$: initial tape-2 positioning, load, tape-0
positioning, reverse copy, and tape-2 return. The loading seed has speed
$1/2$ during the first third and then speed~1, up to bounded raster error.
At final core exit, tape~1 is at most $4L+6K$; tape~0 is at its final
base and tape~2 at zero. Since $K\le L/6$, positioning and output cost
at most $6L$. Root work is at most $11L$.
The clock uses $a=80$. At fixed $r$, its $2^r$ teeth are finite, every
other loop was covered above, and allocation offsets are geometric sums
plus $O(K)$ endpoint terms. Division by $L$ makes the latter vanish.
This proves regularity of every fixed stage and both root stages.

\subsection{C4: disjoint zones, including all transits}

Let $X,Y$ be the fresh tape-0 and tape-1 allocation cursors at a forward
stage, and $Z$ its tape-2 frame. The source base is $X-2M-2$.
The complete stage ranges have lengths at most $4M+3$, $4M+6$ and
$3M+4$. At most seven positionings, and at most $6M$ active events,
give
\[
 |F_r|\le34M+42.
\]
In $B_r$, the tape-0 and tape-1 ranges have lengths at most $6M+5$
and $6M+8$. Tape~2 starts at $Z$ at the deepest stage and at
$Z+(7/2)M+6$ otherwise; its entire range is covered by $6M+8$ too.
There are at most eight positionings. Its comb has exactly $2^r$ teeth
of length $M/2^r$, with $3M-2\cdot2^r$ ticks in total. Initial combine,
both doublings, reverse, merge and copy-back together with the comb
cost at most $14M$. Therefore
\[
 |B_r|\le62M+64,\qquad |F_r|+|B_r|\le96M+106\le128M.
\]
These bounds include cyclic wraps, the moving side heads and their
parked tails. No preservation annotation is used.

Put $S_M=L-2db$. The seed base is $3S_M+4K+2$; complete preparation
has duration $9L-12db+8K+6\le10L$. Before the final positioning,
tape~0 is already at the final base, tape~1 is at $7S_M+10K-3$, and
tape~2 is at $L/2+2$. Positioning plus the $3L$-tick output sweep
costs $10L-14db+10K-3\le10L$. Root work is at most $20L$.

At fixed depth, the comb has a fixed finite number $2^r$ of teeth.
Each tooth has three monotone legs. Positionings are finite unit/stay
pieces, ordinary loops are affine, doubling has bounded raster error,
and the variable loop satisfies \eqref{eq:variable-limit}. All normalized
allocation offsets and shifts converge: they are geometric sums,
$\mathrm{shift}_r/M_r\to1/3$, and endpoint errors $O(K)/L\to0$.
The seed load and reverse copy have finitely affine limits, as does
the final sweep (tapes~0 and~1 have speed $1/3$; tape~2 has speed~1
then parks). Thus all actual fixed stages are regular. The unbounded
\emph{total} number of comb teeth does not violate the theorem's local
fixed-stage hypothesis.

\subsection{C5: comb fanout}

Loading costs $L$ and leaves the work heads at $(L-1,0)$. Assign the
following $L-1$-tick reset to root preparation. This changes only clock
ownership. Every forward and backward recursive stage now starts and ends
at $(0,0)$; subsequent initial resets take no ticks. Root preparation
costs $2L-1$, and the final output sweep costs $L$.

Write $H=M/2$, $N=2^r$, and $h=M/N$. Expanding the literal loops and
positioning calls into constant-move blocks and two-tick repeated blocks
gives
\[
 |F_r|=11H-2,\qquad |B_r|=20H-2N-2.
\]
The backward count includes all $N$ comb teeth, their last stationary
event, both reverse half-loops and the final reset. Thus
$|F_r|+|B_r|\le18M$, and root work $3L-1$ is also at most $18L$.
The forward stage uses 13 blocks; the backward stage uses $4N+10$.
These counts, endpoints and the complete clock partition are proved in
Lean. Their identification with the Rust loops is a source audit, not a
formal semantics for Rust.

For each fixed depth, replace each repeated block by the straight path
with the same average velocity and duration. A block has period one or
two, so its sampled path differs from this ramp by at most two cells.
At most $J_r=4\cdot2^r+13$ blocks suffice for either stage, including
zero-duration padding. The normalized error is at most $2J_r/V_r$,
which tends to zero whenever that stage's charged volume diverges.
Normalized block starts and durations lie in $[0,18]$ and slopes in
$[-1,1]$. Compactness of this finite parameter space supplies a uniform
subsequential limit that is a finite sum of affine ramps.
Consequently every actual fixed stage is regular, for every sequence of
positive entry widths. The unbounded number of stages is paid by
the geometric tail.

\subsection{From the bounds to the clock hypothesis}

The stage ranges above follow inductively from the allocation bases and
the actual last position of each preceding loop. They include the deepest
stage and the first forward entry; these are the endpoints at which an
unpaid reset would otherwise be easy to introduce. With the root budgets,
\cref{tab:bounds} proves a complete clock and \eqref{eq:tail}.

For finite concatenation of stage profiles, convergent normalized endpoints
and durations suffice: all head paths are unit-Lipschitz, so perturbing a
junction time by $\epsilon$ changes a position by at most $\epsilon$.
Each raster error tends to zero, and a fixed stage has only finitely many
pieces. The frozen terminal piece is included. Thus the fixed-stage
regularity requirement holds on every diverging sequence of admitted
dimensions and on every further subsequence. Since the schedules do not depend on the input,
this supplies the complete correct-input quantifiers of
\cref{thm:imported}, with all stages good and no exceptional work.
